% ============================================================
%  Chapter 17 — Bargaining Over Composition Order
% ============================================================
\chapter{Bargaining Over Composition Order}
\label{ch:bargaining}

\begin{quote}
\textit{``The definition of the situation is itself the first move in the game.''}\\
--- Victor Turner, \textit{The Ritual Process} (1969)
\end{quote}

\vspace{1em}

\section{Introduction}

Composition in ideatic space is associative but \emph{not} commutative.
In general, $\compose(a, b) \neq \compose(b, a)$: who speaks first shapes
the meaning of the entire exchange.  This fundamental asymmetry is the
engine of power in signalling games.  The ability to set the agenda---to
determine what is discussed first, which frame is laid down before the other
party responds---is, as Victor Turner recognized, a primary form of social
control.\footnote{Turner's analysis of ritual process emphasizes that the
\emph{order} of ritual stages determines the transformation undergone by
participants.  The liminal phase must precede reaggregation; reversing the
order produces a different---often incoherent---ritual.}

In political discourse, this is the problem of \emph{agenda-setting power}.
In legal reasoning, it is the question of which precedent is applied first.
In Hegelian dialectics, it is the difference between thesis-antithesis and
antithesis-thesis: the reversal produces a different synthesis.

This chapter formalizes the bargaining problem that arises from
non-commutativity.  We introduce the \texttt{compositionOrder} pair, the
\texttt{orderInvariant} predicate, and the \texttt{OrderPair} structure,
then prove eighteen theorems characterizing when order matters, when it
doesn't, and what structural properties hold regardless of ordering.

The key insight is that while the \emph{content} of composition depends on
order, certain \emph{structural} properties---depth bounds, self-resonance,
and the behavior of void---are order-invariant.  This creates a formal
analogy to fair bargaining: the pie may be divided differently depending on
who moves first, but the total size of the pie is bounded the same way
regardless.

\subsection{Axiomatic Bargaining and the Composition Problem}

The mathematical study of bargaining begins with John Nash's landmark
1950 paper, which showed that under four axioms---Pareto efficiency,
symmetry, invariance to affine transformations, and independence of
irrelevant alternatives---there exists a unique bargaining solution, the
\emph{Nash bargaining solution}.\footnote{John F.\ Nash, ``The Bargaining
Problem,'' \textit{Econometrica} 18, no.\ 2 (1950): 155--162.}  Nash's
framework assumes that the bargaining set is fixed and symmetric.  But in
ideatic composition, the bargaining problem is more fundamental: the parties
must agree not merely on a division of surplus but on the \emph{order in
which their contributions are composed}.  Because $\compose(a,b) \neq
\compose(b,a)$ in general, the choice of ordering is itself a strategic
variable, not a parameter.

Ariel Rubinstein's alternating-offers model (1982) extended Nash's static
framework into a dynamic setting, showing that when parties take turns
making proposals, the equilibrium outcome converges to the Nash solution as
the discount factor approaches unity.\footnote{Ariel Rubinstein, ``Perfect
Equilibrium in a Bargaining Model,'' \textit{Econometrica} 50, no.\ 1
(1982): 97--109.}  In our framework, Rubinstein's alternating offers
correspond directly to alternating composition orders: if $A$ proposes
$\compose(a,b)$ and $B$ proposes $\compose(b,a)$, the bargaining process is
a negotiation over which composition is realized.  The key distinction is
that in Rubinstein's model, the proposals converge to a common outcome; in
non-commutative composition, there may be no convergence---the two orderings
may yield genuinely different ideas.

\subsection{Agenda-Setting, Power, and Framing}

Political scientists have long recognized that controlling the
\emph{agenda}---the sequence in which issues are considered---is a primary
form of political power.  William Riker's theory of
\emph{heresthetics}\footnote{William H.\ Riker, \textit{The Art of
Political Manipulation} (1986).} identifies agenda manipulation as one of
three fundamental political strategies: you win not by persuading voters
to change their preferences, but by changing which alternatives they
consider and in what order.  John Kingdon's theory of ``policy
windows''\footnote{John W.\ Kingdon, \textit{Agendas, Alternatives, and
Public Policies}, 2nd ed.\ (2003).} similarly emphasizes that the timing
and sequencing of political proposals determines which policies are adopted.

Steven Lukes's three-dimensional theory of power provides a deeper
framework.\footnote{Steven Lukes, \textit{Power: A Radical View}, 2nd ed.\
(2005).}  The \emph{first dimension} is overt decision-making power; the
\emph{second} is the power to set the agenda, determining which issues are
even considered; the \emph{third} is the power to shape desires and
perceptions, so that dominated parties do not even recognize the
alternatives foreclosed by the agenda.  In IDT terms, the first dimension
corresponds to choosing $a$ versus $b$; the second to choosing
$\compose(a,b)$ versus $\compose(b,a)$; and the third to the ability to
make agents unaware that $\compose(b,a)$ is even a possibility.

Kahneman and Tversky's prospect theory\footnote{Daniel Kahneman and Amos
Tversky, ``Prospect Theory: An Analysis of Decision Under Risk,''
\textit{Econometrica} 47, no.\ 2 (1979): 263--291.} demonstrated that the
\emph{framing} of a decision problem---whether outcomes are presented as
gains or losses relative to a reference point---systematically alters
choices.  Framing is a form of composition ordering: to frame a policy as
``preventing losses'' rather than ``forgoing gains'' is to compose the
evaluative frame \emph{before} the policy content, producing
$\compose(\text{loss-frame}, \text{policy})$ rather than
$\compose(\text{gain-frame}, \text{policy})$.  Our formalism shows that
these two compositions are, in general, genuinely different ideas---not
merely psychologically different presentations of the same idea.

The concept of the ``Overton window''---the range of policies considered
politically acceptable at a given time---can also be understood as a
constraint on composition order.  The window determines which ideas can
serve as the \emph{first} component in a policy composition: only ideas
within the window are available as opening frames.  Shifting the Overton
window is thus equivalent to expanding the set of permissible first-movers
in the bargaining over composition order.

\subsection{Historical Parallels}

The Treaty of Westphalia (1648) offers a paradigmatic illustration.  The
Catholic and Protestant powers negotiated not merely over territorial
boundaries but over which \emph{framework} would serve as the basis for
the emerging state system.  The principle \textit{cuius regio, eius
religio} was itself a meta-bargain: the sovereign's religious identity
would be composed first, and the subjects' beliefs would follow.  In modern
terms, the Westphalian settlement resolved a composition-order dispute by
establishing a fixed ordering rule.

Similarly, the Camp David Accords (1978) required Anwar Sadat and Menachem
Begin to negotiate not just territorial concessions but the \emph{sequence}
in which issues would be addressed.  Jimmy Carter's genius as mediator lay
in restructuring the agenda: by separating the Sinai question from the
Palestinian question and addressing them in a specific order, he changed
the composition of the negotiation itself.\footnote{William B.\ Quandt,
\textit{Camp David: Peacemaking and Politics} (1986).}

In legislative settings, the power to determine the order of business---the
``rules of the game''---is often more consequential than any individual
vote.  The Speaker of the House, the Senate Majority Leader, and committee
chairs all derive power primarily from their ability to set the
\emph{sequence} in which proposals are composed.  Our formalization gives
this intuition mathematical content: the composition order determines the
outcome, and the agenda-setter controls the composition order.

\section{Core Definitions}
\label{sec:ch17-defs}

\begin{definition}[Composition Order Pair]
\label{def:compositionOrder}
The \emph{composition order pair} of $a, b \in \IS$ is
\[
  \texttt{compositionOrder}(a, b) \;:=\;
  \bigl(\compose(a, b),\; \compose(b, a)\bigr) \;\in\; \IS \times \IS.
\]
This captures both possible sequencings of a bilateral exchange: ``I go
first'' versus ``you go first.''
\end{definition}

\begin{lstlisting}[caption={Composition order pair in Lean}]
def compositionOrder (a b : I) : I * I :=
  (IdeaticSpace.compose a b, IdeaticSpace.compose b a)
\end{lstlisting}

\begin{definition}[Order Invariance]
\label{def:orderInvariant}
Two ideas $a, b$ are \emph{order-invariant} when composition commutes:
\[
  \texttt{orderInvariant}(a, b) \;:\Longleftrightarrow\;
  \compose(a, b) = \compose(b, a).
\]
This is a strong condition: most idea-compositions are \emph{not}
order-invariant.  When it holds, there is no agenda-setting advantage.
\end{definition}

\begin{lstlisting}[caption={Order invariance in Lean}]
def orderInvariant (a b : I) : Prop :=
  IdeaticSpace.compose a b = IdeaticSpace.compose b a
\end{lstlisting}

\begin{definition}[Order Pair Structure]
\label{def:OrderPair}
An \emph{order pair} packages two ideas with both their compositions
and proofs that the compositions are correctly formed:
\[
  \texttt{OrderPair} = \bigl(l, r, f, s \;\mid\;
    f = \compose(l, r),\; s = \compose(r, l)\bigr).
\]
\end{definition}

\begin{lstlisting}[caption={Order pair structure in Lean}]
structure OrderPair (I : Type*) [IdeaticSpace I] where
  left : I
  right : I
  first : I
  second : I
  ab_eq : first = IdeaticSpace.compose left right
  ba_eq : second = IdeaticSpace.compose right left
\end{lstlisting}

% ============================================================
\section{Illustrative Examples}
\label{sec:ch17-examples}

Before proceeding to the theorems, we present several extended examples that
illustrate how composition-order bargaining manifests across diverse human
institutions.

\begin{example}[The Treaty of Westphalia]
\label{ex:westphalia}
The Peace of Westphalia (1648) illustrates how composition order shapes
political outcomes. The Catholic and Protestant powers each sought to
establish their theological framework as the ``first mover'' in the
emerging international order. The principle \textit{cuius regio, eius
religio}---the ruler determines the religion---was itself a bargain over
composition order: the sovereign's religious identity is composed first,
and the subjects' beliefs follow. In IDT terms, the sovereign's idea $a$
is composed before the population's idea $b$, yielding $\compose(a, b)$
rather than $\compose(b, a)$---a fundamentally different political
theology.\footnote{The Westphalian system is analyzed through a
game-theoretic lens in Daniel Philpott, \textit{Revolutions in
Sovereignty} (2001).}

Had the composition been reversed---the subjects' beliefs composed before
the sovereign's---the result would have been a form of popular religious
sovereignty, akin to congregationalism.  The entire trajectory of European
political development turned on which ordering was institutionalized.
The Westphalian bargain was, at its core, a bargain over $\compose(a,b)$
versus $\compose(b,a)$.
\end{example}

\begin{example}[Robert's Rules of Order]
\label{ex:roberts_rules}
Parliamentary procedure, codified in Henry M.\ Robert's
\textit{Rules of Order} (1876), is a systematic framework for resolving
composition-order disputes in deliberative assemblies.  The rules specify
an elaborate protocol for determining which proposals are considered first:
the ``order of business'' determines the composition sequence.

A \emph{motion to amend} creates a composition-order problem: should the
original motion $a$ be composed with the amendment $b$ to yield
$\compose(a, b)$, or should the amended version $\compose(b, a)$ be
considered as the base text?  Robert's Rules resolve this by stipulating
that amendments are composed \emph{after} the main motion---the main
motion provides the frame, and the amendment modifies it.  The
\emph{motion to substitute} reverses this ordering entirely: it proposes
to replace $a$ with $b$, effectively switching from $\compose(a, b)$ to
$\compose(b, a)$.

The procedural complexity of parliamentary law---subsidiary motions,
privileged motions, incidental motions---can be understood as a
sophisticated mechanism for managing the exponential growth of possible
composition orderings as the number of proposals increases.  With $n$
proposals, there are $n!$ possible orderings, and the rules of order
serve to reduce this combinatorial explosion to a manageable sequence.
\end{example}

\begin{example}[Media Agenda-Setting: McCombs and Shaw]
\label{ex:media_agenda}
Maxwell McCombs and Donald Shaw's agenda-setting theory (1972)\footnote{Maxwell
E.\ McCombs and Donald L.\ Shaw, ``The Agenda-Setting Function of Mass
Media,'' \textit{Public Opinion Quarterly} 36, no.\ 2 (1972): 176--187.}
demonstrated that the media do not tell people \emph{what} to think, but
rather \emph{what to think about}.  In IDT terms, the media's power lies
not in choosing $a$ over $b$, but in choosing $\compose(a, b)$ over
$\compose(b, a)$: by presenting issue $a$ before issue $b$, the media
establish $a$ as the interpretive frame through which $b$ is subsequently
understood.

Consider a news cycle in which economic data $a$ is presented before a
political scandal $b$.  The public receives $\compose(a, b)$: the scandal
is interpreted through an economic lens.  If the order is reversed---the
scandal is presented first, followed by the economic data---the public
receives $\compose(b, a)$: the economy is interpreted through the lens of
political corruption.  These are, in general, genuinely different public
ideas, not merely different presentations of the same information.

The rise of algorithmic news feeds has intensified this phenomenon.  Each
user's feed presents information in a personalized order, producing
personalized compositions $\compose(a_i, b_i)$ that may differ
dramatically across individuals.  The ``filter bubble'' is, formally, a
personalized composition ordering.
\end{example}

\begin{example}[Socratic Dialogue as Composition-Order Bargaining]
\label{ex:socratic}
In the Platonic dialogues, Socrates consistently positions himself as the
\emph{questioner}, compelling his interlocutors to state their theses
first.  This is a strategic choice of composition order: by ensuring that
the interlocutor's idea $b$ is composed before Socrates's own idea $a$,
Socrates obtains $\compose(a, b)$---his philosophical framework applied
\emph{to} the interlocutor's claim---rather than $\compose(b, a)$, which
would be the interlocutor's framework applied to Socrates's claim.

The \emph{elenchus}---Socrates's method of cross-examination---is a
systematic technique for ensuring that the composition order favors the
questioner.  By asking questions, Socrates forces his interlocutor to
provide the raw material (the first component), which Socrates then
composes with his own dialectical frame (the second component).  The result
is that the interlocutor's thesis is always interpreted \emph{within}
Socrates's framework, never the reverse.

This observation has implications for modern pedagogical theory.  The
``Socratic method'' widely used in legal education is effective precisely
because it enforces a composition order in which the student's initial
response is composed with the professor's analytical frame, producing a
synthesis that the student experiences as self-generated discovery.
\end{example}

% ============================================================
\section{Void and Order Invariance}
\label{sec:ch17-void}

Void---silence, non-contribution---is the universal commuter.  It has no
agenda-setting power because it contributes nothing.

\begin{theorem}[Void is Order-Invariant---Left (17.1)]
\label{thm:void_order_invariant}
For all $a \in \IS$, $\texttt{orderInvariant}(\void, a)$.
\end{theorem}

\noindent\textit{Proof sketch.}\quad
$\compose(\void, a) = a = \compose(a, \void)$ by the identity laws.
\hfill$\square$

\medskip

\noindent\textbf{Commentary.}
Silence has no agenda-setting power.  Whether you speak into silence or
silence follows your speech, the result is the same---just your speech.
In political theory, a party that brings nothing to the table cannot
influence the agenda.

\begin{theorem}[Self-Composition is Order-Invariant (17.2)]
\label{thm:self_order_invariant}
For all $a \in \IS$, $\texttt{orderInvariant}(a, a)$.
\end{theorem}

\noindent\textit{Proof sketch.}\quad
$\compose(a, a) = \compose(a, a)$ trivially.
\hfill$\square$

\medskip

\noindent\textbf{Commentary.}
Internal dialogue with yourself has no ordering problem---you always agree
with yourself on what comes first.  This is the psychological ground truth
that anchors the theory: self-composition is always order-invariant.

\begin{theorem}[Order Invariance of Void Pair (17.8)]
\label{thm:void_void_invariant}
$\texttt{orderInvariant}(\void, \void)$.
\end{theorem}

\noindent\textit{Proof sketch.}\quad
Immediate from both identity laws. \hfill$\square$

\medskip

\noindent\textbf{Commentary.}
Two silences are the same regardless of which comes first.  The degenerate
case anchors the extremal behavior of the theory.

\begin{theorem}[Void is Universally Order-Invariant (17.12)]
\label{thm:void_universal_order_invariant}
For all $a \in \IS$, $\texttt{orderInvariant}(a, \void)$.
\end{theorem}

\noindent\textit{Proof sketch.}\quad
$\compose(a, \void) = a = \compose(\void, a)$.
\hfill$\square$

\medskip

\noindent\textbf{Commentary.}
The absence of contribution never creates ordering conflicts.  Void is
universally transparent to reordering, a formal counterpart to the
observation that silence is the one ``move'' in a game that is costless
in either position.\footnote{In mechanism design, a ``null action'' is
typically assumed to be costless and position-independent. Void formalizes
this assumption within IDT.}

\begin{theorem}[Void Annihilates Order Difference (17.18)]
\label{thm:void_annihilates_order}
For all $a \in \IS$,
\[
  (\texttt{compositionOrder}(a, \void))_1 =
  (\texttt{compositionOrder}(a, \void))_2.
\]
\end{theorem}

\noindent\textit{Proof sketch.}\quad
Both components reduce to $a$ by the identity laws.
\hfill$\square$

\medskip

\noindent\textbf{Commentary.}
A player who contributes nothing has no bargaining power over ordering---there
is nothing to order.  The composition order pair collapses to a diagonal
element when one party contributes void.

% ============================================================
\section{Structural Order Theorems}
\label{sec:ch17-structural}

\begin{theorem}[Order Invariance is Symmetric (17.11)]
\label{thm:order_invariant_symm}
If $\texttt{orderInvariant}(a, b)$, then $\texttt{orderInvariant}(b, a)$.
\end{theorem}

\noindent\textit{Proof sketch.}\quad
If $\compose(a,b) = \compose(b,a)$, then symmetry of equality gives
$\compose(b,a) = \compose(a,b)$.
\hfill$\square$

\medskip

\noindent\textbf{Commentary.}
If swapping $A$ and $B$'s contributions doesn't matter, then swapping $B$
and $A$'s doesn't matter either.  This ensures that order invariance is a
symmetric relation---a necessary condition for it to serve as a formal model
of ``fair exchange.''

\begin{theorem}[Composition Order First Component (17.4)]
\label{thm:order_fst}
$(\texttt{compositionOrder}(a, b))_1 = \compose(a, b)$.
\end{theorem}

\begin{theorem}[Composition Order Second Component (17.5)]
\label{thm:order_snd}
$(\texttt{compositionOrder}(a, b))_2 = \compose(b, a)$.
\end{theorem}

\noindent\textit{Proof sketch.}\quad
Both are definitional equalities.
\hfill$\square$

\medskip

\noindent\textbf{Commentary.}
These projection lemmas are technical necessities, extracting the ``$a$-goes-first''
and ``$b$-goes-first'' orderings from the pair.  They establish the interface
between the order pair abstraction and the underlying composition operation.

\begin{theorem}[Composition Order Swap is Involution (17.17)]
\label{thm:order_swap_involution}
\[
  \texttt{compositionOrder}(b, a) =
  \bigl((\texttt{compositionOrder}(a, b))_2,\;
  (\texttt{compositionOrder}(a, b))_1\bigr).
\]
\end{theorem}

\noindent\textit{Proof sketch.}\quad
Expanding definitions: $\texttt{compositionOrder}(b,a) = (\compose(b,a),
\compose(a,b))$, which equals the swap of $(\compose(a,b), \compose(b,a))$.
\hfill$\square$

\medskip

\noindent\textbf{Commentary.}
The dialectical reversal of a reversal returns to the thesis.  Swapping who
goes first twice recovers the original ordering.  This involution property
ensures that the bargaining structure has no higher-order ordering ambiguities.

\begin{theorem}[Order Pair Construction (17.14)]
\label{thm:order_pair_exists}
For any $a, b \in \IS$, there exists an \texttt{OrderPair} with
$\text{left} = a$ and $\text{right} = b$.
\end{theorem}

\noindent\textit{Proof sketch.}\quad
Construct the pair with $\text{first} = \compose(a,b)$,
$\text{second} = \compose(b,a)$, and the proofs are \texttt{rfl}.
\hfill$\square$

\medskip

\noindent\textbf{Commentary.}
Every bilateral exchange can be formally analyzed by examining both possible
orderings.  The \texttt{OrderPair} structure packages this analysis into a
single object, facilitating the systematic study of agenda-setting effects.

% ============================================================
\section{Power, Framing, and the First-Mover Advantage}
\label{sec:ch17-power}

The preceding structural theorems reveal a striking dichotomy in
composition-order bargaining: the \emph{existence} of an order pair is
guaranteed for any two ideas, yet the \emph{content} of its two components
may diverge arbitrarily.  This section explores the consequences of this
dichotomy for theories of power, framing, and strategic interaction.

\subsection{The First-Mover Advantage in Formal Terms}

In game theory, the \emph{first-mover advantage} refers to the strategic
benefit accruing to the player who acts first in a sequential game.  In
composition-order bargaining, the first-mover advantage is the difference
between $\compose(a, b)$ and $\compose(b, a)$.  When $a$ ``goes first,''
the resulting idea $\compose(a, b)$ is \emph{$a$-framed}: $a$ provides the
interpretive scaffolding onto which $b$ is composed.

The magnitude of the first-mover advantage can be measured by the
``distance'' between the two orderings.  While IDT does not assume a metric
on $\IS$, the resonance relation provides a qualitative measure: when
$\text{resonates}(\compose(a,b), \compose(b,a))$, the first-mover advantage
is ``small'' in the sense that both orderings are at least mutually
coherent.  When the two orderings fail to resonate, the first-mover
advantage is ``large''---the outcomes are so different as to be
incoherent with each other.

Theorem~\ref{thm:resonant_orders} establishes that when the input ideas
resonate, the outputs always resonate.  This provides a formal foundation
for the folk intuition that bargaining between like-minded parties is
``easier'' than bargaining between adversaries: when the parties share a
common frame (resonance), the composition order matters less.

\subsection{Framing as Composition-Order Selection}

The framing literature in behavioral economics identifies two types of
framing effects: \emph{valence framing} (gain vs.\ loss frames) and
\emph{emphasis framing} (which aspects of an issue are highlighted).  Both
can be modeled as composition-order choices.

In valence framing, the same objective outcome $x$ is presented either as
$\compose(\text{gain-frame}, x)$ or $\compose(\text{loss-frame}, x)$.
Kahneman and Tversky's experiments show that these compositions are treated
as different ideas by human cognizers---formally vindicating our axiom that
composition is non-commutative.

In emphasis framing, the question is not which frame to apply but
\emph{when} to apply it.  Presenting a policy's costs before its benefits
yields $\compose(\text{cost}, \text{benefit})$; the reverse yields
$\compose(\text{benefit}, \text{cost})$.  Empirical research consistently
shows that information presented first has disproportionate influence on
final judgments---a primacy effect that is precisely the content-level
manifestation of composition non-commutativity.

\begin{remark}[Connection to Mechanism Design]
\label{rem:mechanism-design}
In mechanism design and implementation theory, the ``designer'' chooses
a game form (mechanism) to implement a desired social choice function.
Composition-order bargaining can be viewed as a special case: the
mechanism designer selects an ordering rule $\sigma: \{a,b\} \to
\{(a,b), (b,a)\}$ that determines which composition is realized.
Theorems~\ref{thm:void_order_invariant}--\ref{thm:void_annihilates_order}
show that void is the unique idea for which the mechanism designer's
choice is irrelevant.  For all other ideas, the designer's ordering
rule has non-trivial consequences---a result reminiscent of the
Gibbard--Satterthwaite theorem's impossibility of
strategy-proofness.\footnote{The connection to Gibbard--Satterthwaite
is suggestive rather than formal.  A full investigation of
composition-order mechanism design is left to future work.}
\end{remark}

\begin{remark}[The Impossibility of Order-Free Composition]
\label{rem:order-free-impossibility}
One might hope for a ``canonical'' composition that is independent of
order---perhaps the set $\{\compose(a,b), \compose(b,a)\}$, or some
average of the two.  But this approach fails for a fundamental reason:
the resulting object is no longer an element of ideatic space $\IS$, but
rather a subset or a distribution over $\IS$.  Order-free composition
requires moving to a higher-order structure---the power set $\mathcal{P}(\IS)$
or the space of distributions over $\IS$.  This observation motivates the
social-choice constructions in Chapter~\ref{ch:social-choice}, where
aggregation rules are introduced precisely to combine multiple compositions
into a single shared idea.
\end{remark}

\begin{remark}[Connections to Gift Exchange and Social Choice]
\label{rem:cross-chapter}
Composition-order bargaining connects directly to the gift-exchange
dynamics of Chapter~\ref{ch:gift-exchange} and the social-choice
mechanisms of Chapter~\ref{ch:social-choice}.  In gift exchange, the
sequential nature of reciprocity creates a natural composition ordering:
the first gift $a$ is composed before the reciprocal gift $b$, and the
social meaning of the exchange $\compose(a, b)$ differs from the meaning
of the reverse sequence $\compose(b, a)$.  The obligation to reciprocate
is, formally, a constraint on composition order.  In social choice, the
aggregation of multiple agents' ideas requires resolving the composition
ordering for all pairs simultaneously---a far more complex problem that
generalizes bilateral bargaining to the $n$-party case.
\end{remark}

% ============================================================
\section{Depth and Resonance Under Reordering}
\label{sec:ch17-depth}

While the \emph{content} of a composition depends on order, its
\emph{depth bounds} are symmetric: subadditivity holds for both orderings.

\begin{theorem}[Both Orders Have Same Depth Bound (17.3)]
\label{thm:both_orders_depth_bound}
For all $a, b \in \IS$,
\begin{align*}
  \depth(\compose(a, b)) &\leq \depth(a) + \depth(b), \\
  \depth(\compose(b, a)) &\leq \depth(a) + \depth(b).
\end{align*}
\end{theorem}

\noindent\textit{Proof sketch.}\quad
The first inequality is subadditivity.  The second follows from
$\depth(\compose(b,a)) \leq \depth(b) + \depth(a) = \depth(a) + \depth(b)$.
\hfill$\square$

\medskip

\noindent\textbf{Commentary.}
The ``cost'' of an exchange, measured by depth, is bounded the same way
regardless of who goes first.  Subadditivity is symmetric.  This means that
while ordering affects the \emph{content} of the outcome, it does not affect
the \emph{resource constraints}---a crucial fairness property.

\begin{theorem}[Depth Bound Under Reordering (17.10)]
\label{thm:depth_reorder_bound}
$\depth(\compose(b, a)) \leq \depth(a) + \depth(b)$.
\end{theorem}

\noindent\textit{Proof sketch.}\quad
From subadditivity with reordering of the bound.
\hfill$\square$

\medskip

\noindent\textbf{Commentary.}
The maximum information capacity of an exchange is invariant under
reordering.  In information-theoretic terms, the channel capacity is
bounded by the sum of the endpoint depths, regardless of direction.

\begin{theorem}[Depth Zero Invariance Under Order (17.16)]
\label{thm:depth_zero_both_orders}
If $\depth(a) = 0$ and $\depth(b) = 0$, then
$\depth(\compose(a,b)) = 0$ and $\depth(\compose(b,a)) = 0$.
\end{theorem}

\noindent\textit{Proof sketch.}\quad
Apply depth-zero closure in both orders.
\hfill$\square$

\medskip

\noindent\textbf{Commentary.}
Trivial exchanges are trivial regardless of who goes first.  You can't
create depth by reordering banalities---a critical-theory insight with
mathematical teeth.

\begin{theorem}[Both Orders Self-Resonate (17.6)]
\label{thm:both_orders_self_resonate}
$\text{resonates}(\compose(a,b), \compose(a,b))$ and
$\text{resonates}(\compose(b,a), \compose(b,a))$.
\end{theorem}

\noindent\textit{Proof sketch.}\quad
Self-resonance is universal.
\hfill$\square$

\medskip

\noindent\textbf{Commentary.}
Any completed exchange achieves internal coherence regardless of who spoke
first.  This is the minimal hermeneutic property: every outcome is
self-consistent.

\begin{theorem}[Resonant Inputs Give Resonant Orders (17.7)]
\label{thm:resonant_orders}
If $\text{resonates}(a, b)$, then
$\text{resonates}\bigl(\compose(a, b),\; \compose(b, a)\bigr)$.
\end{theorem}

\noindent\textit{Proof sketch.}\quad
From \texttt{resonance\_compose} applied to the pairs $(a,b)$ and
$(b,a)$, using the symmetry of resonance.
\hfill$\square$

\medskip

\noindent\textbf{Commentary.}
This is Turner's ritual-process theorem.  When ritual actors are in
resonance---in Turner's terminology, when they share \emph{communitas}---reordering
the ritual stages produces ceremonies that still ``rhyme'' with each other.
The spirit of the ritual is preserved even if the sequence changes.
Formally, resonant inputs produce resonant orderings, ensuring that the
bargaining problem between culturally aligned agents has a coherent
solution space.\footnote{Turner introduced the concept of \emph{communitas}
to describe the intense fellow-feeling that arises in liminal ritual stages.
Our resonance relation formalizes this as a structural property of ideatic space.}

\medskip

\noindent\textbf{Extended discussion of Theorem~\ref{thm:resonant_orders}.}\quad
The resonance-preservation result is the most consequential theorem in this
section for social theory.  It establishes that cultural alignment between
agents is \emph{robust} to composition-order variation: if two agents are
``on the same wavelength'' (resonant), then the two possible orderings of
their exchange are themselves on the same wavelength.  This has three
important implications.

First, it provides a formal criterion for when bargaining over composition
order is \emph{low-stakes}: if the parties resonate, both orderings produce
coherent outcomes, so the distributional consequences of order choice are
bounded.  This aligns with the empirical observation that negotiations
between allies (resonant parties) focus on substance rather than procedure.

Second, it explains why \emph{communitas}---the feeling of deep social
connection---facilitates ritual flexibility.  Turner observed that
communities experiencing strong communitas can improvise ritual sequences
without losing the ritual's transformative power.  Theorem~\ref{thm:resonant_orders}
provides the algebraic explanation: resonance between ritual elements
ensures that reordering preserves coherence.

Third, the contrapositive is equally revealing: if two orderings
$\compose(a,b)$ and $\compose(b,a)$ fail to resonate, then the inputs
$a$ and $b$ must themselves fail to resonate.  Incoherent bargaining
outcomes diagnose incoherent input ideas---a formal version of the
negotiator's maxim that ``if the process breaks down, the parties were
never really aligned.''

% ============================================================
\section{Associativity and Multi-Party Bargaining}
\label{sec:ch17-multi}

\begin{theorem}[Associativity Resolves Three-Party Brackets (17.9)]
\label{thm:three_party_bracket_invariance}
$\compose(\compose(a,b), c) = \compose(a, \compose(b,c))$.
\end{theorem}

\noindent\textit{Proof sketch.}\quad
The associativity axiom of ideatic composition.
\hfill$\square$

\medskip

\noindent\textbf{Commentary.}
In three-party negotiation, whether $(A,B)$ form a coalition first and then
bargain with $C$, or $A$ bargains with a pre-formed $(B,C)$ coalition, the
final composite meaning is identical.  Associativity is the great equalizer:
it eliminates the ``coalition formation order'' as a source of bargaining
advantage, even as the \emph{within-pair} ordering remains significant.

\begin{theorem}[Triple Depth Bound---Both Bracketings (17.13)]
\label{thm:triple_depth_bound}
$\depth(\compose(\compose(a,b), c)) \leq \depth(a) + \depth(b) + \depth(c)$.
\end{theorem}

\noindent\textit{Proof sketch.}\quad
Cascade two applications of subadditivity.
\hfill$\square$

\medskip

\noindent\textbf{Commentary.}
Three-way collaboration has bounded complexity no matter how sub-teams are
formed.  The depth bound generalizes linearly, providing a universal
resource constraint on multi-party cooperation.

\begin{theorem}[Resonance Between Reordered Triples (17.15)]
\label{thm:triple_reorder_resonance}
If $\text{resonates}(a, b)$ and $\text{resonates}(b, c)$, then
\[
  \text{resonates}\bigl(\compose(\compose(a,b), c),\;
  \compose(a, \compose(b,c))\bigr).
\]
\end{theorem}

\noindent\textit{Proof sketch.}\quad
By associativity, both sides are equal, so self-resonance applies.
\hfill$\square$

\medskip

\noindent\textbf{Commentary.}
When all ritual participants share communitas, rearranging the ritual order
produces ceremonies that deeply resonate.  The resonance between re-bracketed
triples is not merely an analogy to Turner's theory---it is its formal
content.

% ============================================================

\noindent\textbf{Extended discussion of Theorem~\ref{thm:three_party_bracket_invariance}.}\quad
The associativity result for three-party bracketing has deep implications
for coalition theory.  In real-world multi-party negotiations---such as
the formation of coalition governments in parliamentary democracies---the
question of \emph{which sub-coalitions form first} is a major strategic
concern.  Associativity tells us that while the \emph{within-pair} ordering
matters (who speaks first within a dyad), the \emph{between-pair} bracketing
does not: whether Labor and the Greens form a coalition first and then
negotiate with the Liberals, or Labor negotiates with a pre-formed
Green-Liberal alliance, the final composite position is the same.

This result provides formal support for the observation in coalition theory
that the ``order of coalition formation'' matters less than the ``order of
proposal within each coalition.''  The bracketing structure drops out;
the sequencing within each bilateral exchange does not.  This is a
non-trivial constraint on the strategy space of multi-party bargaining:
parties should focus on their bilateral ordering advantages rather than
on the meta-strategic question of which coalitions form when.

The depth bound for triples (Theorem~\ref{thm:triple_depth_bound})
extends this analysis to resource constraints: the total complexity of a
three-party exchange is bounded by the sum of individual complexities,
regardless of bracketing.  This linearity is essential for scalability:
it means that the resource cost of multi-party bargaining grows
additively, not multiplicatively, ensuring that large-scale negotiations
remain tractable.

\section{Conclusion}
\label{sec:ch17-concl}

This chapter has established eighteen theorems characterizing the
bargaining problem that arises from the non-commutativity of ideatic
composition.

The central tension is clear: while \emph{content} depends on order
(who speaks first matters), \emph{structural properties}---depth bounds,
self-resonance, void transparency---do not.  This creates a formal
separation between the distributional question (``who gets the better
outcome?'') and the efficiency question (``how much total meaning is
produced?'').

The key results are:
\begin{itemize}
  \item Void has no agenda-setting power
    (Theorems~\ref{thm:void_order_invariant}--\ref{thm:void_annihilates_order}).
  \item Depth bounds are order-invariant
    (Theorems~\ref{thm:both_orders_depth_bound}--\ref{thm:depth_zero_both_orders}).
  \item Resonant agents produce resonant orderings
    (Theorem~\ref{thm:resonant_orders}).
  \item Associativity resolves multi-party bracketing
    (Theorem~\ref{thm:three_party_bracket_invariance}).
\end{itemize}

Turner's ritual process reveals that transformative social experiences
require a specific ordering of stages.  Our formalization shows that this
ordering dependence is a fundamental feature of ideatic composition, not
an artifact of particular rituals.  In the next chapter, we turn from
bilateral bargaining to the problem of diffusion through networks.
